
% ============================================================================
% CHAPTER 1: SYSTEM ARCHITECTURE
% ============================================================================
\chapter{SYSTEM ARCHITECTURE}
\label{ch:architecture}

\section{Overview}

The Multi-Agent Data Analysis and Insight Generation Platform operates as a sequential pipeline of six specialized agents. Each agent processes data through a shared state object called GraphState, which maintains all intermediate results and metadata. This design allows agents to operate independently while sharing computed facts, enabling modularity and fault isolation.

\section{Data Flow and Shared State}

Agent 1 begins with raw CSV data. It profiles column distributions, detects implicit missing values encoded as sentinels, and identifies column relationships. The structural profile and cached DataFrame pass to Agent 2. Agent 2 infers semantic meaning for each column using heuristics and optional LLM consultation, producing a schema blueprint with confidence scores. Agent 3 cleans the data according to the schema: it handles missing values, detects and clips outliers, derives business metrics, and logs every transformation. Agent 4 analyzes the cleaned data statistically and generates chart specifications. Agent 5 validates data integrity against business rules and computes quality metrics. Agent 6 combines all prior results into a final report with LLM-generated narrative and visual insights.

All agents operate deterministically except where LLM calls introduce stochasticity. A streaming progress API allows real-time monitoring of pipeline execution.

% ============================================================================
% CHAPTER 2: AGENT IMPLEMENTATIONS
% ============================================================================
\chapter{AGENT IMPLEMENTATIONS}
\label{ch:agents}

\section{Agent 1: Data Profiling}

Agent 1 ingests raw CSV files and produces a structural understanding of the dataset.

\subsection{Distribution Analysis}

Agent 1 classifies numeric columns by computing skewness and kurtosis. Right-skewed columns have positive skewness greater than 1. Normal columns satisfy both $|\text{skewness}| \leq 0.5$ and $|\text{kurtosis}| \leq 1$. All others are symmetric. This classification drives downstream outlier-clipping strategy: percentile bounds for skewed data, IQR $1.5\times$ for normal data.

\subsection{Implicit Missingness Detection}

Some datasets encode missing values as special numbers or strings. Numeric columns are checked against common sentinels: -999, -9999, -99, -1, 999, and others. Text fields match patterns like ``n/a'', ``none'', ``missing'', ``#n/a''. Agent 1 flags these per-column with affected row counts. This catches quality issues that simple null checks miss, especially in financial data where -999 is standard practice.

\subsection{Column Relationships}

Agent 1 identifies potential primary keys (entirely unique columns), duplicate column pairs (identical data under different names), and strong correlations ($|r| \geq 0.5$). These relationships reveal schema structure and redundant columns.

\section{Agent 2: Schema and LLM Reliability}

Agent 2 infers semantic meaning for each column and handles LLM communication reliably.

\subsection{LLM Failover Strategy}

Agent 2 tries Groq API first, then Gemini. For Gemini, users supply multiple API keys via environment variables: \texttt{GEMINI\_API\_KEYS} (comma/space/semicolon-separated list), or \texttt{GEMINI\_API\_KEY\_1} through \texttt{GEMINI\_API\_KEY\_5}. When one key hits quota, the system rotates to the next without retry backoff. This avoids HTTP 429 rate-limit cascades.

\subsection{Schema Confidence Scores}

Each semantic tag receives a confidence score (0--100) based on whether the LLM tag matches the heuristic type detector, accounting for dtype, sample values, and missing rates. Low-confidence tags flag columns needing human review.

\subsection{LLM Batching for Wide Schemas}

Datasets with 100+ columns are split into batches to avoid token-limit truncation. If an LLM response truncates or fails to parse, the batch size halves recursively. This handles wide datasets without manual column exclusion.

\section{Agent 3: Data Cleaning}

Agent 3 cleans and transforms the data according to the schema.

\subsection{Adaptive Outlier Clipping}

Clipping strategy depends on distribution shape. Right-skewed columns use percentile bounds (5th to 95th). Normal columns use $1.5\times$ IQR. This avoids over-clipping rare-but-real values in skewed data while still catching true outliers. Bounds are logged for audit.

\subsection{Business Metrics Derivation}

When revenue, cost, quantity, or price columns are detected, Agent 3 automatically computes profit (revenue $-$ cost), margin (profit $\div$ revenue), discount (listed-price $-$ sale-price), and per-unit revenue. Column name matching uses whole-word logic to avoid false matches.

\subsection{Preprocessing Profiles}

Users select strictness at runtime: strict (tighter thresholds, max currency 1.1B), balanced (default), or lenient (looser thresholds, max currency 10B). The system detects whether the dataset is financial or generic and suggests a profile, but users override as needed. Reconciliation tolerance also changes per profile (1\%, 2\%, 5\%).

\subsection{Column Ledger}

Every transformation logs before/after null counts, parse failure rates, range-check failures, clipping bounds, and notes. This audit trail shows where data quality degrades and makes results reproducible.

\subsection{Row Safety Checks}

Deduplication and filtering include safeguards. If a single step removes more than 50\% of rows, the pipeline stops with a diagnostic message, preventing silent data loss from bugs.

\section{Agent 4: Statistical Analysis and Chart Planning}

Agent 4 analyzes the cleaned data and generates visualization specifications.

\subsection{Statistical Chart Selection}

Agent 4 does not use keyword lists. Instead, it scores eight chart families on statistical strength:

\begin{itemize}
	\item \textbf{Dimension ranking}: ANOVA $\eta^2$ (category groups vs. numeric metric)
	\item \textbf{Pareto}: top-N value concentration
	\item \textbf{Distribution}: histogram (spread, skew, outliers)
	\item \textbf{Trend}: regression $R^2$ (metric over time)
	\item \textbf{Seasonality}: coefficient of variation (month effects)
	\item \textbf{Correlation scatter}: Pearson $r$ (strongest pairs)
	\item \textbf{Crosstab heatmap}: Cramér's $V$ (two category columns)
	\item \textbf{Anomaly overlay}: statistical outliers
\end{itemize}

Builders are independent, so failures in one family do not block others. Charts are ranked and capped at a target count (e.g., 12 charts per report).

\subsection{Materiality Gating}

A ``clear leader'' claim requires either $1.5\times$ higher value or $5+$ percentage points ahead of second place. Pareto charts group weak performers ($<5\%$ share) as ``Other''. This filters out trivial differences.

\subsection{Trend Wording}

Even statistically significant trends with $<3\%$ actual change are labeled ``flat''. Statistical significance does not always mean business relevance.

\subsection{Unified Chart Specifications}

All families output the same ChartSpec JSON format. Data is pre-aggregated so static PNG and interactive ECharts renderers consume identical bytes, preventing inconsistencies between views.

\section{Agent 5: Validation and Data Quality}

Agent 5 validates data against business rules and computes quality metrics.

\subsection{Temporal Logic Validation}

Impossible date sequences flag data defects: orders before creation dates, shipments before orders, deliveries before shipments.

\subsection{Formula Reconciliation}

The system checks quantity $\times$ price $\approx$ revenue within 2\% mean absolute error. Mismatches suggest missing cost data or column artifacts.

\subsection{Near-Perfect Correlation Handling}

Correlations $> 0.98$ usually indicate duplicates or data quality signals, not business relationships. They are logged to quality metrics but excluded from narrative findings.

\subsection{Broken Rule Filtering}

Validation rules firing on $>95\%$ of rows are marked as likely misconfigured. They remain in audit logs but are removed from issue counts, preventing garbage rules from drowning out real problems.

\subsection{Cohen's Kappa}

The system computes Cohen's kappa between LLM semantic tags and the local heuristic type detector. Low agreement flags columns needing review.

\section{Agent 6: Report Generation}

Agent 6 combines all prior results into a final report.

\subsection{Narrative Claim Grounding}

Before publishing, Agent 6 extracts numeric claims from the LLM narrative (e.g., ``revenue increased 25\%'') and verifies them against deterministic facts. Flagged hallucinations are logged with examples.

\subsection{Contradiction Detection}

The narrative is scanned for opposite claims. High-contradiction counts flag LLM inconsistencies.

\subsection{Chart Deduplication}

Identical chart specs appearing in multiple sections render only once, removing redundancy and reducing report size.

\subsection{Dynamic Glossary}

Instead of a full glossary, only terms used in that report are defined. Maximum 12 entries, matched against narrative text.

\subsection{Section Ordering}

Sections reorder by combined chart priority (strongest signal first). Readers encounter impactful findings upfront rather than following fixed order.

% ============================================================================
% CHAPTER 3: DATA PROCESSING AND QUALITY
% ============================================================================
\chapter{DATA PROCESSING AND QUALITY}
\label{ch:processing}

\section{Data Ingestion}

The platform ingests CSV and Excel files with multiple resilience strategies.

\subsection{Mixed Delimiters}

The CSV reader handles rows with inconsistent delimiters (comma vs. semicolon). When field counts mismatch, the reader attempts to reconstruct rows using both delimiters.

\subsection{Multi-Encoding Fallback}

CSV ingestion tries UTF-8, cp1252, and Latin-1 in sequence. This handles international business datasets.

\section{Quality Metrics and Validation}

A composite Data Quality Score (0--100) combines remaining null percentage, validation failure percentage, and duplicate rate, weighted by preprocessing profile (strict, balanced, or lenient). Reliability metadata include per-stage confidence scores and decision readiness flags (ready, needs-review, blocked).

\section{Audit Trails}

Every transformation logs before/after null counts, parse failures, range-check failures, imputation methods, scaling bounds, and notes. The ColumnLedger enables reproducibility and identifies where quality degrades through the pipeline.

% ============================================================================
% CHAPTER 4: APIs AND SERVICES
% ============================================================================
\chapter{APIs AND SERVICES}
\label{ch:apis}

\section{RESTful Endpoints}

\subsection{Analysis Pipeline}

\begin{description}
	\item[POST /api/analyze] Upload CSV or Excel and optional preprocessing profile; returns job ID.
	\item[GET /api/analyze/\{job\_id\}/stream] Server-Sent-Events progress updates: agent completions, charts ready, validation status, report generation.
	\item[GET /api/analyze/\{job\_id\}/result] Fetch full analysis results including statistics, charts, quality, and reliability metadata.
\end{description}

\subsection{Job Management}

\begin{description}
	\item[GET /api/jobs] List all jobs (completed and in-progress).
	\item[GET /api/jobs/\{job\_id\}] Fetch single job metadata.
	\item[DELETE /api/jobs/\{job\_id\}] Delete job and artifacts.
	\item[DELETE /api/jobs] Delete all jobs.
\end{description}

\subsection{Report and Artifacts}

\begin{description}
	\item[GET /api/report/\{job\_id\}] Download HTML or PDF report (PDF falls back to HTML if conversion fails).
	\item[GET /api/plots/\{file\_path\}] Serve static chart PNG (per-job isolation under \texttt{outputs/charts/\{run\_id\}/}).
\end{description}

\subsection{Infrastructure}

\begin{description}
	\item[GET /api/health] LLM provider health check (Groq, Gemini, HuggingFace, PostgreSQL); results cached 90 seconds.
	\item[POST /api/health/test-llm] Manual provider health test (bypasses cache).
\end{description}

\section{Dataset Q\&A Service}

After analysis completes, users ask questions about the data. The chat service builds topic-scoped context from relevant sections (correlation, trend, ranking) based on detected question topics, then queries the LLM with grounded facts to prevent hallucination.

\subsection{On-Demand Charts}

Users can request new charts within the chat service. Allowed types: bar, line, histogram, box, scatter. Allowed operations: count, sum, mean, median, min, max, nunique. Whitelists prevent unexpected requests and data abuse.

\section{Progress Streaming}

Server-Sent-Events on \texttt{/api/analyze/\{job\_id\}/stream} provide real-time pipeline progress without polling. Events include agent progress, chart generation, validation status, and report readiness.

% ============================================================================
% CHAPTER 5: INFRASTRUCTURE AND CONFIGURATION
% ============================================================================
\chapter{INFRASTRUCTURE AND CONFIGURATION}
\label{ch:infrastructure}

\section{LLM Provider Management}

\subsection{Multi-Model Failover}

The system attempts Groq API first, then Gemini. This provides redundancy if one provider is unavailable or rate-limited.

\subsection{Multi-Gemini-Key Rotation}

Environment variables \texttt{GEMINI\_API\_KEYS} (comma/space/semicolon-separated list) or \texttt{GEMINI\_API\_KEY\_1} through \texttt{GEMINI\_API\_KEY\_5} enable key rotation without code changes. When one key hits quota, the system advances to the next without retry backoff, preventing 429 rate-limit cascades.

\subsection{Gemini Model Fallback Chain}

The system tries \texttt{gemini-flash-latest} first, then \texttt{gemini-2.5-flash}, \texttt{gemini-2.5-flash-lite}, and \texttt{gemini-2.0-flash}. This gracefully handles model deprecation.

\subsection{Health Check Caching}

Provider health checks cache results for 90 seconds to avoid quota exhaustion on free-tier accounts. Manual checks available via \texttt{POST /api/health/test-llm}.

\section{Optional Persistence}

Job metadata, authentication, and chat history can store in PostgreSQL (via \texttt{DATABASE\_URL} environment variable) or remain in-memory. This supports multi-instance deployments.

\section{Semantic Search}

Row embeddings and analysis facts store in PostgreSQL pgvector. Two document families are indexed: dataset rows and analysis facts (summaries, correlations, trends, anomalies). This enables semantic dataset discovery across jobs.

\section{Preprocessing Profiles}

Users select preprocessing strictness at runtime via \texttt{analysis\_config.preprocessing\_profile}. This propagates through the entire pipeline.

\section{LLM Context Scoping}

The chat service attaches only relevant analysis sections based on detected question topics. This reduces token usage and grounds responses in dataset facts.
